\chapter{核心代码实现}\label{app:code}

本附录给出三大 PINN 模型的关键代码片段，供读者复现参考。完整代码见项目仓库 \texttt{src/models/}。

% ─────────────────────────────────────────────────────────────────────────────
\section{PDE 残差计算（以 BSM 为例）}\label{app:pde_residual}

以下代码展示了如何通过 PyTorch 自动微分精确计算 BSM PDE 残差，无需手动推导有限差分格式：

\begin{verbatim}
def _pde_residual(self, S, t):
    S.requires_grad_(True)
    t.requires_grad_(True)
    # 归一化输入
    x = torch.cat([S / self.S_max, t / self.T], dim=1)
    V = self.net(x)

    # 一阶导数：dV/dt, dV/dS
    V_t = torch.autograd.grad(
        V, t, grad_outputs=torch.ones_like(V),
        create_graph=True)[0]
    V_S = torch.autograd.grad(
        V, S, grad_outputs=torch.ones_like(V),
        create_graph=True)[0]
    # 二阶导数：d2V/dS2
    V_SS = torch.autograd.grad(
        V_S, S, grad_outputs=torch.ones_like(V_S),
        create_graph=True)[0]

    # BSM PDE 残差
    residual = (V_t
        + 0.5 * self.sigma**2 * S**2 * V_SS
        + self.r * S * V_S
        - self.r * V)
    return residual
\end{verbatim}

关键点：\texttt{create\_graph=True} 使得高阶导数可以继续参与反向传播，从而在训练时对 PDE 残差损失求梯度。\texttt{requires\_grad\_(True)} 将 $S$ 和 $t$ 标记为需要求导的叶节点。

% ─────────────────────────────────────────────────────────────────────────────
\section{ICPINN 硬约束输出变换（Heston）}\label{app:icpinn}

以下代码展示了 Heston 模型的硬约束输出变换，确保 $\hat{V}(S,v,T) = \phi(S)$ 精确成立：

\begin{verbatim}
class HestonNet(nn.Module):
    def __init__(self, K, T, hidden=128, depth=6):
        super().__init__()
        self.K = K
        self.T = T
        layers = [nn.Linear(3, hidden), nn.Tanh()]
        for _ in range(depth - 1):
            layers += [nn.Linear(hidden, hidden), nn.Tanh()]
        layers.append(nn.Linear(hidden, 1))
        self.net = nn.Sequential(*layers)

    def forward(self, S_norm, v_norm, t_norm, S_raw, T):
        x = torch.cat([S_norm, v_norm, t_norm], dim=1)
        net_out = self.net(x)
        # 终值收益：payoff = max(S - K, 0)
        payoff = torch.clamp(S_raw - self.K, min=0.0)
        # 硬约束：当 t=T 时 (T-t)=0，输出精确等于 payoff
        tau = T - t_norm * T   # 剩余时间
        return payoff + tau * net_out
\end{verbatim}

与软约束方式相比，硬约束无需在损失函数中加入 $\mathcal{L}_\text{ic}$ 项，减少了损失函数中各项之间的梯度冲突，收敛更稳定。

% ─────────────────────────────────────────────────────────────────────────────
\section{LLM 路由层系统提示词}\label{app:llm_prompt}

以下是 LLM 路由层使用的系统提示词（system prompt），用于约束 DeepSeek 模型的输出格式：

\begin{verbatim}
You are an option pricing assistant. Extract parameters from
the user's description and output ONLY valid JSON.

Model selection rules:
- If user provides only constant volatility sigma: use "BSM"
- If user mentions beta or "CEV": use "CEV"
- If user provides kappa/theta/xi/rho/v0 or mentions
  "Heston" or "stochastic volatility": use "Heston"
- Default: "BSM"

Required JSON fields:
{
  "model": "BSM" | "CEV" | "Heston",
  "option_type": "call" | "put",
  "S": float,   // current stock price
  "K": float,   // strike price
  "T": float,   // time to expiry in years
  "r": float,   // risk-free rate
  "sigma": float,          // BSM/CEV only
  "beta": float,           // CEV only, default 0.5
  "kappa": float,          // Heston only
  "theta": float,          // Heston only
  "xi": float,             // Heston only
  "rho": float,            // Heston only
  "v0": float              // Heston only
}
\end{verbatim}

通过 \texttt{response\_format: \{type: "json\_object"\}} 参数约束输出格式，确保 LLM 始终返回可解析的 JSON，而非自然语言描述。

% ─────────────────────────────────────────────────────────────────────────────
\section{美式期权 PINN 扩展方案}\label{app:american}

美式期权的定价问题是一个自由边界问题，期权价值满足变分不等式：
\begin{equation}
  \min\left(\mathcal{F}[V],\; V - \phi(S)\right) = 0,
\end{equation}
其中 $\phi(S) = \max(S-K, 0)$ 为立即行权价值，$\mathcal{F}$ 为 BSM 算子。

在 PINN 框架中，可通过在损失函数中加入互补条件惩罚项处理这一约束：
\begin{equation}
  \mathcal{L}_\text{american} = \mathcal{L}_\text{pde}
    + \lambda \cdot \frac{1}{N_c}\sum_{i=1}^{N_c}
      \left[\min\left(\mathcal{F}[\hat{V}](x_i),\;
      \hat{V}(x_i) - \phi(x_i)\right)\right]^2,
\end{equation}
其中 $\lambda$ 为惩罚权重。该方法由 Louskos\citep{louskos2021american} 验证，可在不修改网络架构的情况下处理提前行权约束。
